% concatenated from HdR-curve.tex
\documentclass[12pt,reqno]{amsart}
\usepackage{hdrcurve}
\title[Hodge-to-de Rham degeneration and quasihomogeneity]{Hodge-to-de Rham degeneration and quasihomogeneous singularities of curves}
\author{Yunfan He}
\address{Beijing International Center for Mathematical Research, Peking University}
\email{he\_yunfan@pku.edu.cn}

\begin{document}
\begin{abstract}
We study the Hodge-to-de Rham spectral sequence for integral projective curves with local complete intersection singularities. We prove that degeneration at the \(E_2\)-page is equivalent to requiring every singularity to be a quasihomogeneous plane curve singularity. We also show that, in the same local complete intersection setting, the Hochschild-to-cyclic spectral sequence degenerates at the \(E_2\)-page if and only if the same condition holds.
\end{abstract}
\maketitle

\section{Introduction}
The study of algebraic de Rham cohomology originates with Grothendieck. In \cite[Theorem~1']{Grothendieck_66}, he proved that for a smooth scheme \(X\) over \(\mathbb{C}\), the hypercohomology
\[
H^{\bullet}(X, \Omega_X^{\bullet})
\]
of the complex of sheaves of differentials \(\Omega_X^{\bullet}\) computes the singular cohomology of the analytification \(X^{\mathrm{an}}\). This complex is called the \emph{algebraic de Rham complex}, and carries the stupid/Hodge filtration, which induces a spectral sequence
\[
E_1^{p,q}(X) = H^q(X, \Omega_X^p) \Longrightarrow H^{p+q}_{\mathrm{sing}}(X^{\mathrm{an}},\mathbb{C}),
\]
known as the \emph{Hodge-to-de Rham} spectral sequence. Grothendieck \cite[p.9]{Grothendieck_66} further showed that if \(X\) is smooth and projective over \(\mathbb{C}\), then this spectral sequence degenerates at the first page. Deligne and Illusie \cite[Corollary~2.7]{Deligne-Illusie_87} later extended this degeneracy result to smooth proper schemes over any field of characteristic \(0\), using reduction to positive characteristic.

For singular varieties, the correct replacement for the algebraic de Rham complex is given by the Hodge-completed derived de Rham complex \(\widehat{\mathrm{dR}}_X^{\bullet}\) \cite[\S 4.1]{Bhatt_12}. Its Hodge filtration has graded pieces
\[
\operatorname{gr}_F^p \widehat{\mathrm{dR}}_X^{\bullet}\simeq \bigwedge^{p}\mathbb{L}^{\bullet}_{X/\mathbb{C}}[-p],
\]
and therefore yields the (derived) Hodge-to-de Rham spectral sequence
\[
E_1^{p,q}(X) = H^q\!\left(X,\bigwedge^p\mathbb{L}^{\bullet}_X\right) \Longrightarrow H^{p+q}_{\mathrm{sing}}(X^{\mathrm{an}},\mathbb{C}).
\]
He \cite[Theorem~1.1]{He_25} proved that nodal projective curves have \(E_2\)-degeneration, and the same phenomenon appears for the cuspidal cubic curve; see \cite[Theorem~7.1]{He_25}. It is therefore natural to ask for a singularity-theoretic characterization of degeneration at the second page. The following theorem gives a complete answer in the local complete intersection case.

\begin{theorem}
For an integral projective curve \(X\) with local complete intersection singularities, its Hodge-to-de Rham spectral sequence degenerates at the \(E_2\)-page if and only if every singularity of \(X\) is a quasihomogeneous plane curve singularity.
\end{theorem}

Besides the Hodge filtration on derived de Rham cohomology, one can also consider the filtration on negative cyclic homology arising from the Hochschild mixed complex of \(X\). This yields the \emph{Hochschild-to-cyclic} spectral sequence. In the same local complete intersection setting, we show that degeneration at the \(E_2\)-page is characterized by exactly the same singularity-theoretic condition.

\begin{corollary}
For an integral projective curve \(X\) with local complete intersection singularities, the following are equivalent:
\begin{enumerate}
\item the Hodge-to-de Rham spectral sequence of \(X\) degenerates at the \(E_2\)-page;
\item the Hochschild-to-cyclic spectral sequence of \(X\) degenerates at the \(E_2\)-page;
\item every singularity of \(X\) is a quasihomogeneous plane curve singularity.
\end{enumerate}
\end{corollary}

The proof of the main theorem contains two parts. In the plane case, we describe the negative-row tail differentials on the \(E_1\)-page in terms of the Milnor and Tjurina algebras of the singularities. For a quasihomogeneous plane curve singularity, the weighted Euler relation identifies these local tail maps with multiplication by nonzero scalars on weighted pieces of the Milnor algebra, so they are isomorphisms. Conversely, if the spectral sequence degenerates at the \(E_2\)-page, a global dimension count forces Milnor number equal to the Tjurina number at every planar singularity, hence quasihomogeneity. In the non-planar local complete intersection case, we show that a singularity of maximal embedding dimension produces a nonzero \(E_2\)-term in total degree \(>2\), which is impossible for a projective curve. 

Section~\ref{sec:singularities} reviews the singularity-theoretic background used in the proof, with emphasis on plane curve singularities and quasihomogeneity. Section~\ref{sec:e2-degeneration} studies the Hodge-to-de Rham spectral sequence. In Section~\ref{subsec:plane-case} we analyze the \(E_1\)-page and its \(d_1\)-differentials in the plane case. In Section~\ref{subsec:non-planar-case} we show that a non-planar local complete intersection singularity forces a nonzero \(E_2\)-term in total degree \(>2\). In Section~\ref{subsec:main-theorem} we combine these analyses to prove the main theorem, and in Section~\ref{subsec:HC} we study the Hochschild-to-cyclic spectral sequence and prove the companion corollary. Section~\ref{sec:further-directions} records several questions suggested by these results.


%\subsection*{Acknowledgements}
%We would like to thank


\section{Singularities on integral projective curves}\label{sec:singularities}

In this section we collect the classes of curve singularities that will be used in the proof of the main theorem.

\subsection{Basic definitions}

Let $X$ be an integral projective curve over $\mathbb{C}$, and let $x\in X$ be a singular point. 

\begin{definition}
We say that $x$ is a \emph{plane curve singularity} if its completed local ring is of the form
\[
\widehat{\mathcal O}_{X,x}\cong \mathbb{C}[[u,v]]/(f)
\]
for some nonzero power series $f\in \mathbb{C}[[u,v]]$.
\end{definition}

\begin{definition}
A power series
\[
f(u,v)=\sum c_{ab}u^a v^b
\]
is called \emph{weighted homogeneous} if there exist positive rational numbers $w_u,w_v$ and \(d\) such that every monomial $u^av^b$ appearing in $f$ satisfies
\[
aw_u+bw_v=d.
\]
\end{definition}

\begin{remark}
If \(f\) is weighted homogeneous of weighted degree \(d\) with weights \(w_u,w_v\), then after dividing the weights by \(d\) one obtains positive rational numbers \(w_1,w_2\) such that the \emph{Euler relation}
\begin{equation}\label{eq:weighted-euler}
f=w_1u f_u+w_2v f_v
\end{equation}
holds. This relation plays a central role in the analysis of the Hodge-to-de Rham spectral sequence.
\end{remark}

\begin{definition}
A plane curve singularity is called \emph{quasihomogeneous} if, after a formal change of coordinates, it can be defined by a weighted homogeneous equation.
\end{definition}

\begin{definition}
We say that $x\in X$ is a \emph{local complete intersection singularity} if $\mathcal O_{X,x}$ is a local complete intersection ring.
\end{definition}

\begin{remark}
If $x$ is a plane curve singularity, then $x$ is a hypersurface singularity, hence a local complete intersection singularity.
However, the converse is false in general.
\end{remark}

\subsection{Examples}

Among the standard examples of quasihomogeneous plane curve singularities are the simple ADE singularities; see, for example, \cite[\S~10]{Milnor_68}:
\[
A_n:\ u^2+v^{n+1},
\qquad
D_n:\ u^2v+v^{n-1},
\qquad
E_6:\ u^3+v^4,
\]
\[
E_7:\ u^3+uv^3,
\qquad
E_8:\ u^3+v^5.
\]
Each of these is weighted homogeneous, hence quasihomogeneous and planar.

Thus one has the chain of inclusions of types of singularities
\[
\{\text{simple}\}
\subset
\{\text{quasihomogeneous plane}\}
\subset
\{\text{plane}\}
\subset
\{\text{lci}\}.
\]

\begin{example}
Let
\[
X:=\Spec \mathbb{C}[[t^3,t^4,t^5]],
\]
and let \(x\in X\) be its unique closed point. Then
\[
\widehat{\mathcal O}_{X,x}\cong \mathbb{C}[[t^3,t^4,t^5]]
\]
has embedding dimension \(3\). Hence \(x\) is not a plane curve singularity.
\end{example}

\subsection{Local numerical invariants}

Let \(x\in X\) be a singular point, and let \(\widetilde{\mathcal O}_{X,x}\) denote the normalization of \(\mathcal O_{X,x}\). We define

\begin{itemize}
  \item the delta-invariant by
    \[
    \delta_x:=\dim_{\mathbb{C}}\!\left(\widetilde{\mathcal O}_{X,x}/\mathcal O_{X,x}\right),
    \]
  \item the number of branches \(r_x\) to be the number of points of the normalization lying over \(x\),
\end{itemize}

Now assume \((X,x)\) is a plane curve singularity such that
\[
\widehat{\mathcal O}_{X,x}\cong \mathbb{C}[[u,v]]/(f)
\]
with \(f\) reduced; we refer to such a singularity as a \emph{reduced plane curve singularity}. We further define

\begin{itemize}
  \item the Milnor algebra:
    \[
    M_f:=\mathbb{C}[[u,v]]/(f_u,f_v),
    \]
  \item the Milnor number:
    \[
    \mu_x:=\dim_{\mathbb{C}}M_f,
    \]
  \item the Tjurina algebra:
    \[
    T_f:=\mathbb{C}[[u,v]]/(f,f_u,f_v)\cong M_f/fM_f,
    \]
  \item the Tjurina number:
    \[
    \tau_x:=\dim_{\mathbb{C}}T_f.
    \]
\end{itemize}

For plane curve singularities, Milnor proved in his celebrated book \cite[\S~10]{Milnor_68} the famous Milnor's formula:
\begin{equation}\label{eq:milnor-formula}
\mu_x=2\delta_x-r_x+1.
\end{equation}

Moreover, Saito \cite{Saito_71} proved the following criterion for quasihomogeneity of reduced plane curve singularities:
\begin{equation}\label{eq:tau-mu-qh}
\tau_x=\mu_x
\quad\Longleftrightarrow\quad
x \text{ is quasihomogeneous}.
\end{equation}

\subsection{Global setup}

Let $X$ be an integral projective curve, and let
\[
\nu:\widetilde X\to X
\]
be its normalization. We write
\[
g:=g(\widetilde X).
\]
If $X$ has singular locus
\[
\mathrm{Sing}(X)=\{x_1,\dots,x_s\},
\]
then the arithmetic genus of \(X\) is 
\[
p_a(X)=g+\sum_{i=1}^s \delta_{x_i}.
\]
See, for example, \cite[Chapter~IV]{Hartshorne_77}.

\section{\texorpdfstring{$E_2$}{E2}-degeneration for integral projective lci curves}\label{sec:e2-degeneration}

In this section we prove the main result of the paper.

Recall that for a smooth projective variety \(W\) of dimension \(n\), the \emph{de Rham complex} is the complex of sheaves
\[
0 \to \mathcal O_W \to \Omega_W^1 \to \Omega_W^2 \to \cdots \to \Omega_W^n \to 0,
\]
where \(\Omega_W^i\) denotes the sheaf of K\"ahler \(i\)-forms on \(W\).

For a singular variety \(X\), the analogous object is the \emph{Hodge-completed derived de Rham complex}
\[
\widehat{\mathrm{dR}}_X^{\bullet} : 0 \to \mathcal O_X \to \mathbb{L}^{\bullet}_X \to \bigwedge^2 \mathbb{L}^{\bullet}_X \to \bigwedge^3 \mathbb{L}^{\bullet}_X \to \cdots.
\]
The Hodge filtration on \( \widehat{\mathrm{dR}}_X^{\bullet} \) induces a spectral sequence whose \(E_1\)-page is
\[
E_1^{p,q}(X)
   = H^{q}\!\left(X, \bigwedge^{p}\mathbb{L}^{\bullet}_X\right),
\]
and which converges to the singular cohomology \(H^{p+q}_{\mathrm{sing}}(X^{\mathrm{an}},\mathbb C)\) (see \cite[\S~5]{Bhatt_12}). However, no completion is required for the integral projective lci curves \(X\), since the cotangent complex \(\BLb_X\) is bounded. We refer to this as the \emph{Hodge-to-de Rham spectral sequence}, and write \(E_r^{p,q}(X)\) for its \(r\)-th page.

Our goal is to determine when this spectral sequence degenerates at the \(E_2\)-page. In Section~\ref{subsec:plane-case} we analyze the \(E_1\)-page and its \(d_1\)-differentials in the plane case. In Section~\ref{subsec:non-planar-case} we show that a non-planar local complete intersection singularity forces a nonzero \(E_2\)-term in total degree \(>2\). We then combine these two analyses to prove the main theorem in Section~\ref{subsec:main-theorem}, and finally in Section~\ref{subsec:HC} we compare the Hodge-to-de Rham and Hochschild-to-cyclic spectral sequences.


\subsection{The plane case}\label{subsec:plane-case}

In this section, let \(X\) be an integral projective curve over \(\mathbb{C}\) whose singularities are plane curve singularities. Recall we use \(\mathbb{L}^{\bullet}_X\) to denote the cotangent complex of \(X\), and \(\Omega_X^1:=\mathcal{H}^0(\mathbb{L}^{\bullet}_X)\) to denote the cotangent sheaf of \(X\). For each singular point \(x\in \mathrm{Sing}(X)\), we denote
\[
A=A_x:=\widehat{\mathcal O}_{X,x}\cong R/(f),
\qquad
R=\mathbb C[[u,v]],
\]
and, when no confusion, we suppress the dependence on \(x\). We also use the notation \(M_f\) and \(T_f\) from Section~2.3 for the corresponding Milnor and Tjurina algebras.

We begin by recording the basic global properties of the cotangent complex itself.

\begin{lemma}\label{lem:plane-cotangent}
The cotangent complex \(\mathbb{L}^{\bullet}_X\) of \(X\) is perfect of amplitude \([-1,0]\), and
\[
\mathcal H^{-1}(\mathbb{L}^{\bullet}_X)=0.
\]
\end{lemma}

\begin{proof}
Since every singular point of \(X\) is planar and \(X\) is smooth elsewhere, \(X\) is lci. Hence the cotangent complex \(\mathbb{L}^{\bullet}_X\) is perfect of amplitude \([-1,0]\). It remains to show that \(\mathcal H^{-1}(\mathbb{L}^{\bullet}_X)=0\).

On the smooth locus of \(X\), one has \(\mathbb{L}^{\bullet}_X\simeq \Omega_X^1[0]\), so there is nothing to prove. Let \(x\in \mathrm{Sing}(X)\). Since \(\mathcal O_{X,x}\) is a reduced local ring essentially of finite type over \(\mathbb C\), it is excellent, hence Nagata \cite[\href{https://stacks.math.columbia.edu/tag/07QV}{Lemma 07QV}]{stacks-project}. Therefore its completion \(A=\widehat{\mathcal O}_{X,x}\) is reduced \cite[\href{https://stacks.math.columbia.edu/tag/07NZ}{Lemma 07NZ}]{stacks-project}. Thus in the presentation \(A\cong \mathbb{C}[[u,v]]/(f)\), the element \(f\) may be chosen to be reduced.
By the standard hypersurface description of the cotangent complex \cite[\href{https://stacks.math.columbia.edu/tag/08RB}{Lemma 08RB}]{stacks-project},
\begin{equation}\label{eq:plane-hypersurface-model}
\mathbb{L}^{\bullet}_A\simeq \left[A\xrightarrow{(f_u,f_v)}A^{\oplus 2}\right],
\end{equation}
with \(A^{\oplus 2}\) in degree \(0\). We claim that the map \(A\to A^{\oplus 2}\) is injective. 

Indeed, an element \(a\in A\) lies in the kernel exactly when \(af_u=af_v=0\) in \(A\), equivalently when for a lift \(\widetilde a\in \mathbb{C}[[u,v]]\) one has \(\widetilde a f_u,\widetilde a f_v\in (f)\). Since \(f\) is reduced, it is square-free, so no irreducible factor of \(f\) divides both \(f_u\) and \(f_v\); equivalently, \(\gcd(f,f_u,f_v)=1\). Hence every irreducible factor of \(f\) divides \(\widetilde a\), and therefore \(\widetilde a\in (f)\). Thus \(a=0\), proving the claim. It follows that
\[
\mathcal H^{-1}(\mathbb{L}^{\bullet}_A)=0.
\]

Recall that flat base change for the cotangent complex gives
\[
\mathbb{L}^{\bullet}_A \simeq (\mathbb{L}^{\bullet}_X)_x\otimes_{\mathcal O_{X,x}}^{\mathbf L} A.
\]
Since \((\mathbb{L}^{\bullet}_X)_x\) has amplitude \([-1,0]\) and \(A\) is flat over \(\mathcal O_{X,x}\), taking degree \(-1\) cohomology commutes with this base change:
\[
\mathcal H^{-1}(\mathbb{L}^{\bullet}_A)\cong \mathcal H^{-1}(\mathbb{L}^{\bullet}_X)_x\otimes_{\mathcal O_{X,x}} A.
\]
Indeed, one may represent \((\mathbb{L}^{\bullet}_X)_x\) by a two-term complex of finite free \(\mathcal O_{X,x}\)-modules in degrees \([-1,0]\), and then tensoring with the flat \(\mathcal O_{X,x}\)-algebra \(A\) preserves exactness in degree \(-1\).
Since completion is faithfully flat, \(\mathcal H^{-1}(\mathbb{L}^{\bullet}_X)_x=0\). Thus \(\mathcal H^{-1}(\mathbb{L}^{\bullet}_X)=0\).

\end{proof}

We next turn to the higher derived exterior powers, which are responsible for the negative rows of the Hodge-to-de Rham spectral sequence.

\begin{lemma}\label{lem:plane-exterior-powers}
For every \(p\ge 2\), the derived exterior power \(\bigwedge^p\mathbb{L}^{\bullet}_X\) is supported on \(\mathrm{Sing}(X)\). More explicitly, let \(x\in \mathrm{Sing}(X)\), and use the standing notation above.
Then the completed stalk of \(\bigwedge^p\mathbb{L}^{\bullet}_X\) at \(x\) is quasi-isomorphic to the three-term complex
\[
\left[
A \xrightarrow{\binom{f_u}{f_v}} A^{\oplus 2}
\xrightarrow{(-f_v\ \ f_u)} A
\right],
\]
placed in degrees \([-p,-p+2]\). In particular,
\[
H^q\!\left(X,\bigwedge^p\mathbb{L}^{\bullet}_X\right)=0
\qquad\text{for } q\notin\{-p+1,-p+2\}.
\]
\end{lemma}

\begin{proof}
On the smooth locus of \(X\), one has \(\mathbb{L}^{\bullet}_X\simeq \Omega_{X_{\mathrm{sm}}}^1[0]\), so
\[
\bigwedge^p\mathbb{L}^{\bullet}_X\big|_{X_{\mathrm{sm}}}\simeq 0
\qquad (p\ge 2).
\]
Hence \(\bigwedge^p\mathbb{L}^{\bullet}_X\) is supported on \(\mathrm{Sing}(X)\).

As in the proof of Lemma~\ref{lem:plane-cotangent}, flat base change gives
\[
(\mathbb{L}^{\bullet}_X)_x\otimes_{\mathcal O_{X,x}}^{\mathbf L} A \simeq \mathbb{L}^{\bullet}_A.
\]
Since \(\mathbb{L}^{\bullet}_X\) is perfect, taking derived exterior powers commutes with this flat base change, so
\[
\left(\bigwedge^p\mathbb{L}^{\bullet}_X\right)_x\otimes_{\mathcal O_{X,x}}^{\mathbf L} A
\simeq
\bigwedge^p\!\left((\mathbb{L}^{\bullet}_X)_x\otimes_{\mathcal O_{X,x}}^{\mathbf L} A\right).
\]
Therefore the completed stalk of \(\bigwedge^p\mathbb{L}^{\bullet}_X\) at \(x\) is isomorphic to \(\bigwedge^p\mathbb{L}^{\bullet}_A\).

Now using the hypersurface description~\eqref{eq:plane-hypersurface-model}, we can compute the higher wedge powers as 
\[
\bigwedge^p\mathbb{L}^{\bullet}_A \simeq
\left[
A \xrightarrow{\binom{f_u}{f_v}} A^{\oplus 2}
\xrightarrow{(-f_v\ \ f_u)} A
\right],
\]
placed in degrees \([-p,-p+2]\). By the injectivity argument from Lemma~\ref{lem:plane-cotangent}, the leftmost map is injective, so the cohomology of this complex can be nonzero only in degrees \(-p+1\) and \(-p+2\).
For each \(i\), since \(A=\widehat{\mathcal O}_{X,x}\) is flat over \(\mathcal O_{X,x}\), we have
\[
\mathcal H^i\!\left(\bigwedge^p\mathbb{L}^{\bullet}_X\right)_x\otimes_{\mathcal O_{X,x}} A
\cong
H^i\!\left(\bigwedge^p\mathbb{L}^{\bullet}_A\right).
\]
Since completion is faithfully flat, it follows that
\[
\mathcal H^i\!\left(\bigwedge^p\mathbb{L}^{\bullet}_X\right)_x=0
\qquad\text{for } i\notin\{-p+1,-p+2\}.
\]
As \(x\in \mathrm{Sing}(X)\) was arbitrary, the only possibly nonzero cohomology sheaves of \(\bigwedge^p\mathbb{L}^{\bullet}_X\) are \(\mathcal H^{-p+1}\) and \(\mathcal H^{-p+2}\), both supported on the finite set \(\mathrm{Sing}(X)\). Consider the hypercohomology spectral sequence
\[
E_2^{a,b}=H^a\!\left(X,\mathcal H^b\!\left(\bigwedge^p\mathbb{L}^{\bullet}_X\right)\right)
\Longrightarrow
H^{a+b}\!\left(X,\bigwedge^p\mathbb{L}^{\bullet}_X\right).
\]
Since \(H^a(X,\mathcal H^b(\bigwedge^p\mathbb{L}^{\bullet}_X))=0\) for all \(a>0\), this spectral sequence degenerates at \(E_2\), and therefore
\[
H^q\!\left(X,\bigwedge^p\mathbb{L}^{\bullet}_X\right)
\cong
H^0\!\left(X,\mathcal H^q\!\left(\bigwedge^p\mathbb{L}^{\bullet}_X\right)\right).
\]
In particular, \(H^q(X,\bigwedge^p\mathbb{L}^{\bullet}_X)\) can be nonzero only for \(q=-p+1\) or \(q=-p+2\).
\end{proof}

Combining Lemmas~\ref{lem:plane-cotangent} and~\ref{lem:plane-exterior-powers}, we obtain the shape of the \(E_1\)-page.

\begin{corollary}\label{cor:plane-e1-shape}
The \(E_1\)-page of the Hodge-to-de Rham spectral sequence of \(X\) has the form
\[
\begin{tikzcd}[row sep=tiny, column sep=tiny]
1
  & H^1(\cO_X) \ar[r, "d_1^{0,1}"]
  & H^1(\Omega_X^1)
\\[0.2em]
0
  & H^0(\cO_X) \ar[r,"0"]
  & H^0(\Omega_X^1) \ar[r, "d_1^{1,0}"]
  & H^0(X,\bigwedge^2\mathbb{L}^{\bullet}_X)
\\[0.2em]
-1
  &
  &
  &
  H^{-1}(X,\bigwedge^2\mathbb{L}^{\bullet}_X) \ar[r,"d_1^{1,-1}"]
  & H^{-1}(X,\bigwedge^3\mathbb{L}^{\bullet}_X)
\\[0.2em]
-2
  &
  &
  &
  &
  H^{-2}(X,\bigwedge^3\mathbb{L}^{\bullet}_X) \ar[r,"d_1^{2,-2}"]
  & H^{-2}(X,\bigwedge^4\mathbb{L}^{\bullet}_X)
\\[0.2em]
\vdots
  &
  &
  &
  &
  &
  \ddots
\end{tikzcd}
\]
Among the \(d_1\)-differentials on the \(E_1\)-page, only the displayed maps can be nonzero, and all other \(d_1\)-maps vanish for degree reasons.
\end{corollary}

\begin{proof}
By Lemma~\ref{lem:plane-cotangent}, one has \(\mathbb{L}^{\bullet}_X\simeq \Omega_X^1[0]\). Therefore the \(0\)-th and \(1\)-st columns are
\[
E_1^{0,q}(X)=H^q(X,\mathcal O_X),
\qquad
E_1^{1,q}(X)=H^q(X,\Omega_X^1).
\]
By Lemma~\ref{lem:plane-exterior-powers}, for every \(p\ge 2\) the complex \(\bigwedge^p\mathbb{L}^{\bullet}_X\) is supported on the finite set \(\mathrm{Sing}(X)\), so
\[
E_1^{p,q}(X)=0
\qquad\text{unless}\qquad
q=-p+1 \text{ or } q=-p+2.
\]
Since \(X\) is a projective curve, \(H^q(X,\mathcal O_X)=H^q(X,\Omega_X^1)=0\) for \(q\notin\{0,1\}\). Since \(d_1\) has bidegree \((1,0)\), the displayed pattern follows immediately.
\end{proof}

With above description of the \(E_1\)-page, in order to prove the \(E_2\)-degeneartion of the Hodge-to-de Rham spectral sequence, we start by analyzing the infinite tail maps 
\[
d_1^{p,-p}:H^{-p}(X,\bigwedge^{p+1}\mathbb{L}^{\bullet}_X)\to H^{-p}(X,\bigwedge^{p+2}\mathbb{L}^{\bullet}_X).
\]
We will just denote it by \(d_1\) for simplicity. We start by computing these cohomology groups locally at the singular points.

\begin{lemma}\label{lem:tail-map-quotient}
Let
\[
  R=\mathbb C[[u,v]],\qquad A=R/(f),
\]
where $f\in R$ defines an isolated plane curve singularity. We use
\[
  (0:_{M_f}f):=\{m\in M_f\mid fm=0\}
\]
to denote the annihilator of \(f\) in \(M_f\). Then for every $p\ge 1$ there are canonical isomorphisms
\[
H^{-p}\!\left(\bigwedge^{p+1}\mathbb{L}^{\bullet}_A\right)\cong (0:_{M_f}f),
\qquad
H^{-p}\!\left(\bigwedge^{p+2}\mathbb{L}^{\bullet}_A\right)\cong T_f,
\]
In particular, their dimensions both equal to the Tjurina number:
\[
  \dim_{\mathbb C}H^{-p}\!\left(\bigwedge^{p+1}\mathbb{L}^{\bullet}_A\right)
  =
  \dim_{\mathbb C}H^{-p}\!\left(\bigwedge^{p+2}\mathbb{L}^{\bullet}_A\right)
  =
  \tau_x.
\]
\end{lemma}
  
\begin{proof}
We choose the standard dg algebra resolution 
\[
S:=R[\varepsilon],\qquad |\varepsilon|=-1,\qquad \partial(\varepsilon)=f.
\]
Then
\[
\mathbb{L}^{\bullet}_A\simeq \Omega^1_{S}\otimes_S A,
\]
hence can be represented by the two-term complex
\[
A\, d\varepsilon
\xrightarrow{(f_u,f_v)}
A\,du\oplus A\,dv,
\]
in degree \([-1,0]\). For \(m\ge 0\), write
\[
\gamma_m:=\frac{(d\varepsilon)^m}{m!}.
\]
viewed as a normalized element of \(\operatorname{Sym}^m(A\cdot d\varepsilon)\). For each $p\ge 1$, the standard formula for derived exterior powers of this two-term complex gives a three-term complex representing \(\bigwedge^{p+1}\mathbb{L}^{\bullet}_A\), namely
\[
K_p[-p+1],
\]
where
\[
  K_p:
  A\cdot \gamma_{p+1}
  \longrightarrow
  (A\,du\oplus A\,dv)\cdot \gamma_p
  \longrightarrow
  A\cdot du\wedge dv\,\gamma_{p-1},
  \]
  with differentials
  \[
  c\gamma_{p+1}\longmapsto
  (cf_u\,du+cf_v\,dv)\gamma_p,
  \]
  \[
  (a\,du+b\,dv)\gamma_p
  \longmapsto
  (-af_v+bf_u)\,du\wedge dv\,\gamma_{p-1}.
  \]
  After forgetting the harmless basis elements $\gamma_i$, this is just the Koszul complex
\[
K:
A \xrightarrow{\binom{f_u}{f_v}} A^{\oplus 2}\xrightarrow{(-f_v\ \ f_u)} A.
\]

Notice that $K$ is precisely
\[
A\otimes_R K_R(f_u,f_v),
\]
where $K_R(f_u,f_v)$ is the Koszul complex on \(f_u,f_v\). Since the singularity is isolated, the Jacobian ideal \((f_u,f_v)\) is \(\mathfrak m_R\)-primary. As \(R=\mathbb C[[u,v]]\) is a regular local ring of dimension \(2\), the ideal \((f_u,f_v)\) has height \(2\), so \(f_u,f_v\) form a regular sequence in \(R\). Since
\[
M_f=R/(f_u,f_v),
\]
the Koszul complex \(K_R(f_u,f_v)\) is a free resolution of \(M_f\). Therefore, after tensoring with \(A\), it computes the derived tensor product \(A\otimes_R^{\mathbf L}M_f\), so
\[
K \simeq A\otimes_R^{\mathbf L} M_f.
\]
On the other hand, the two-term free resolution
\[
[R\xrightarrow{\cdot f} R]
\]
of $A=R/(f)$ gives
\[
A\otimes_R^{\mathbf L}M_f \simeq [M_f\xrightarrow{\cdot f} M_f],
\]
with degree \(-1\) and \(0\). Therefore, in \(D(A)\),
\[
\bigwedge^{p+1}\mathbb{L}^{\bullet}_A\simeq [M_f\xrightarrow{\cdot f} M_f][-p+1].
\]
Consequently,
\[
  H^{-p}\!\left(\bigwedge^{p+1}\mathbb{L}^{\bullet}_A\right)
  \cong
  \ker\!\bigl(\cdot f:M_f\to M_f\bigr)
  =
  \{m\in M_f\mid fm=0\}
  =(0:_{M_f}f),
  \]
  and
  \[
  H^{-p+1}\!\left(\bigwedge^{p+1}\mathbb{L}^{\bullet}_A\right)
  \cong
  \operatorname{coker}\!\bigl(\cdot f:M_f\to M_f\bigr)
  =M_f/fM_f=T_f.
  \]
  
  Replacing $p$ by $p+1$, we also get
  \[
  \bigwedge^{p+2}\mathbb{L}^{\bullet}_A\simeq [M_f\xrightarrow{\cdot f} M_f][-p],
  \]
  hence
  \[
  H^{-p}\!\left(\bigwedge^{p+2}\mathbb{L}^{\bullet}_A\right)\cong T_f.
  \]
  Since \(M_f\) is finite-dimensional and \(\cdot f:M_f\to M_f\) is an endomorphism, its kernel and cokernel have the same dimension. Since \(T_f\) has dimension \(\tau_x\), the two displayed isomorphisms show that both groups have dimension \(\tau_x\).
  \end{proof}

\begin{lemma}\label{lem:tail-map-euler}
Keep the setting of Lemma~\ref{lem:tail-map-quotient}, and assume further that \(f=w_1u f_u+w_2v f_v\) is weighted homogeneous. Then for every $p\ge 1$, the cohomology groups of Lemma~\ref{lem:tail-map-quotient} are both isomorphic to \(M_f\):
\[
H^{-p}\!\left(\bigwedge^{p+1}\mathbb{L}^{\bullet}_A\right)\cong M_f,
\qquad
H^{-p}\!\left(\bigwedge^{p+2}\mathbb{L}^{\bullet}_A\right)\cong M_f.
\]
Moreover, under these identifications, the de Rham differential induces an isomorphism
\[
d_1:
H^{-p}\!\left(\bigwedge^{p+1}\mathbb{L}^{\bullet}_A\right)
\longrightarrow
H^{-p}\!\left(\bigwedge^{p+2}\mathbb{L}^{\bullet}_A\right).
\]
\end{lemma}

\begin{proof}
By Lemma~\ref{lem:tail-map-quotient},
\[
H^{-p}\!\left(\bigwedge^{p+1}\mathbb{L}^{\bullet}_A\right)\cong (0:_{M_f}f),
\qquad
H^{-p}\!\left(\bigwedge^{p+2}\mathbb{L}^{\bullet}_A\right)\cong T_f.
\]
Since Euler relation~\eqref{eq:weighted-euler} gives \(f\in (f_u,f_v)\), multiplication by \(f\) on \(M_f\) is zero. Therefore
\[
(0:_{M_f}f)=M_f,
\qquad
T_f=M_f/fM_f\cong M_f,
\]
which gives the two displayed identifications in the statement.

To compute the de Rham differential, we use the explicit three-term complex from the proof of Lemma~\ref{lem:tail-map-quotient}:
\[
K:
A \xrightarrow{\Delta_2} A^{\oplus 2}\xrightarrow{\Delta_1} A,
\]
where
\[
\Delta_2(c)=(cf_u,cf_v),\qquad \Delta_1(a,b)=-af_v+bf_u.
\]
There we identified
\[
H^{-p}\!\left(\bigwedge^{p+1}\mathbb{L}^{\bullet}_A\right)=H^{-1}(K),
\qquad
H^{-p}\!\left(\bigwedge^{p+2}\mathbb{L}^{\bullet}_A\right)=H^{0}(K).
\]
By Lemma~\ref{lem:tail-map-quotient} and the Euler relation, these two groups are already abstractly isomorphic to \(M_f\). We now construct an explicit identification \(M_f\cong H^{-1}(K)\), which will allow us to compute the de Rham differential on the chain level.

Next, consider the Euler syzygy
\[
E:=(-w_2v,w_1u)\in A^{\oplus 2}.
\]
Because
\[
\Delta_1(E)=w_2v f_v+w_1u f_u=f=0\quad\text{in }A,
\]
the element $E$ defines a class in $H^{-1}(K)$. Multiplication by $E$ yields a map
\[
\theta:M_f\longrightarrow H^{-1}(K),\qquad [h]\longmapsto [hE].
\]
This map is well-defined, since in \(A^{\oplus 2}\) one has
\[
f_uE=\Delta_2(-w_2v),
\qquad
f_vE=\Delta_2(w_1u).
\]
We claim that \(\theta\) is surjective. Let \([(a,b)]\in H^{-1}(K)\), so
\[
-af_v+bf_u=0
\qquad\text{in }A.
\]
Choose lifts \(\widetilde a,\widetilde b\in R\). Then
\[
-\widetilde a f_v+\widetilde b f_u\in (f),
\]
so there exists \(\widetilde c\in R\) such that
\[
-\widetilde a f_v+\widetilde b f_u=\widetilde c\,f.
\]
Using the Euler relation~\eqref{eq:weighted-euler}, we get
\[
(\widetilde b-\widetilde c\,w_1u)f_u+(-\widetilde a-\widetilde c\,w_2v)f_v=0
\qquad\text{in }R.
\]
As noted in the proof of Lemma~\ref{lem:tail-map-quotient}, \((f_u,f_v)\) is a regular sequence in \(R\). Therefore there exists \(\widetilde d\in R\) such that
\[
\widetilde b-\widetilde c\,w_1u=-\widetilde d\,f_v,
\qquad
-\widetilde a-\widetilde c\,w_2v=\widetilde d\,f_u.
\]
After modulo \((f)\), we obtain
\[
b-\overline{\widetilde c}\,w_1u=-\overline{\widetilde d}\,f_v,
\qquad
-a-\overline{\widetilde c}\,w_2v=\overline{\widetilde d}\,f_u
\qquad
\text{in }
A.
\]
Equivalently,
\[
(a,b)=\overline{\widetilde c}\,E+\Delta_2(-\overline{\widetilde d}),
\]
so \([(a,b)]=\theta([\overline{\widetilde c}])\). Hence \(\theta\) is surjective.

Since \(M_f\) is finite-dimensional and \(H^{-1}(K)\cong (0:_{M_f}f)=M_f\), the surjective map \(\theta:M_f\to H^{-1}(K)\) is an isomorphism, and
\[
H^{-1}(K)\cong M_f.
\]

We now compute the de Rham differential on this explicit model. Since
\[
d_{\mathrm{dR}}(d\varepsilon)=0,
\]
one has \(d_{\mathrm{dR}}(\gamma_m)=0\) for every \(m\ge 0\). Therefore the ordinary de Rham differential induces a chain map
\[
\delta_p:K[-p+1]\longrightarrow K[-p]
\]
whose nontrivial components are given by
\[
\begin{tikzcd}[column sep=large, row sep=large]
A
  \arrow[r, "\Delta_2"]
  \arrow[d, "\delta_p^{-p-1}" left]
&
A^{\oplus 2}
  \arrow[r, "\Delta_1"]
  \arrow[d, "\delta_p^{-p}"]
&
A
  \arrow[d, "0"]
\\
A^{\oplus 2}
  \arrow[r, "\Delta_1"]
&
A
  \arrow[r, "0"]
&
0,
\end{tikzcd}
\]
where
\[
\delta_p^{-p-1}(c):=(c_u,c_v),
\]
\[
\delta_p^{-p}(a,b):=b_u-a_v.
\]
Indeed,
\[
d_{\mathrm{dR}}(c)=c_u\,du+c_v\,dv
\]
gives \(\delta_p^{-p-1}(c)=(c_u,c_v)\), while
\[
d_{\mathrm{dR}}(a\,du+b\,dv)=(b_u-a_v)\,du\wedge dv.
\]
Under the identification \(H^{-1}(K)\cong M_f\) via \(\theta\), the induced map on cohomology is computed by
\[
[h]\longmapsto [\delta_p^{-p}(hE)].
\]
Since
\[
hE=(-h w_2v,\; h w_1u),
\]
we obtain
\[
\delta_p^{-p}(hE)=\partial_u(hw_1u)-\partial_v(-hw_2v)
=
w_1(uh_u+h)+w_2(vh_v+h).
\]
Thus
\[
d_1([h])
=
\bigl[w_1(uh_u+h)+w_2(vh_v+h)\bigr]
\in M_f.
\]

If $[h]$ is weighted homogeneous of weighted degree $\lambda$, then Euler's formula gives
\[
w_1uh_u+w_2vh_v=\lambda h
\quad\text{in }M_f,
\]
hence
\begin{equation}\label{eq:weighted-tail-scalar}
d_1([h])=(\lambda+w_1+w_2)[h].
\end{equation}
Since $f_u$ and $f_v$ are weighted homogeneous, the Jacobian ideal $(f_u,f_v)$ is graded for the induced weighted grading on $R$, so
\[
M_f=\bigoplus_{\lambda}(M_f)_\lambda
\]
is a finite direct sum of weighted-homogeneous pieces. By \eqref{eq:weighted-tail-scalar}, $d_1$ acts on $(M_f)_\lambda$ by the scalar \(\lambda+w_1+w_2\), which is nonzero because \(w_1,w_2>0\) and \(\lambda\ge 0\). Therefore $d_1$ is an isomorphism.
\end{proof}
  
Lemma~\ref{lem:tail-map-euler} gives the local tail isomorphisms in the quasihomogeneous case. We next describe the differential \(d_1^{1,0}\) for an arbitrary plane curve.

Let \(\mathcal A\) and \(\mathcal D\) be the sheaves on \(X\) defined by 
\[
\mathcal A:=\mathcal H^{-1}\!\left(\bigwedge^2\mathbb{L}^{\bullet}_X\right),
\qquad
\mathcal D:=\mathcal H^0\!\left(\bigwedge^2\mathbb{L}^{\bullet}_X\right),
\] and denote by \(\omega_X\) the dualizing sheaf of \(X\).

\begin{lemma}\label{lem:plane-d10-comparison}
Assume that \(X\) has only plane curve singularities. Define the morphism \(\alpha\) by the following composite:
\[
\alpha:\mathcal A
\xrightarrow{\ \mathcal H^{-1}(d_{\mathrm{dR}})\ }
\mathcal H^{-1}\!\left(\bigwedge^3\mathbb{L}^{\bullet}_X\right)
\xrightarrow{\sim}
\mathcal D,
\]
where the second arrow is the canonical isomorphism described in the proof below. Then \(\mathcal A\) and \(\mathcal D\) are supported on \(\mathrm{Sing}(X)\), and there is an exact sequence of sheaves on \(X\)
\[
0\to \mathcal A\to \Omega_X^1\to \omega_X\to \mathcal D\to 0.
\]
Moreover, there is a canonical isomorphism
\[
H^0\!\left(X,\bigwedge^2\mathbb{L}^{\bullet}_X\right)\cong H^0(X,\mathcal D),
\]
under which the restriction of
\[
d_1^{1,0}:H^0(X,\Omega_X^1)\to H^0\!\left(X,\bigwedge^2\mathbb{L}^{\bullet}_X\right)
\]
along the inclusion \(H^0(X,\mathcal A)\hookrightarrow H^0(X,\Omega_X^1)\) is identified with
\[
H^0(\alpha):H^0(X,\mathcal A)\to H^0(X,\mathcal D).
\]
\end{lemma}

\begin{proof}
On the smooth locus of \(X\), one has
\[
\bigwedge^2\mathbb{L}^{\bullet}_X\simeq \Omega_X^2[0]=0,
\]
so \(\mathcal A\) and \(\mathcal D\) are supported on \(\mathrm{Sing}(X)\).
Fix \(x\in \operatorname{Sing}(X)\), and after completion write \(A\cong \mathbb{C}[[u,v]]/(f)\). By \eqref{eq:plane-hypersurface-model}, written with respect to the basis \(du,dv\) of \(A^{\oplus 2}\),
\[
\mathbb{L}^{\bullet}_A\simeq \left[A\xrightarrow{(f_u,f_v)}A\,du\oplus A\,dv\right].
\]
Hence \(\bigwedge^2\mathbb{L}^{\bullet}_A\) is represented by
\[
\left[
A\xrightarrow{(f_u,f_v)}
A\,du\oplus A\,dv
\xrightarrow{(-f_v\ \ f_u)}
A\,du\wedge dv
\right][-1].
\]
Since \(R=\mathbb{C}[[u,v]]\) is regular of dimension \(2\), one can compute the dualizing module of \(A\) as 
\[
\omega_A\cong \operatorname{Ext}^1_R(A,\omega_R)\cong A\,du\wedge dv.
\]
Under this identification, the map
\[
A\,du\oplus A\,dv\longrightarrow A\,du\wedge dv
\]
is given by \(\alpha\mapsto df\wedge \alpha\). Since \(df=f_u\,du+f_v\,dv\), it kills the
submodule \(A\cdot df\subset A\,du\oplus A\,dv\), because \(df\wedge df=0\). Hence it descends to a map
\[
c_A:\Omega_A^1=(A\,du\oplus A\,dv)/(A\cdot df)\longrightarrow \omega_A,
\]
and this is the canonical class map, sending \(a\,du+b\,dv\) to \((bf_u-af_v)\,du\wedge dv\)
(compare \cite[\href{https://stacks.math.columbia.edu/tag/0E9X}{Section 0E9X}]{stacks-project}). Therefore the cohomology sequence of the above three-term complex yields an exact sequence on the completed stalks
\[
0\to \mathcal A_x\to (\Omega_X^1)_x\to \omega_{X,x}\to \mathcal D_x\to 0.
\]
Since completion is faithfully flat, exactness may be checked after tensoring with
\(A=\widehat{\mathcal O}_{X,x}\). Hence the completed exact sequence implies the corresponding exact sequence on the ordinary stalks, and these glue to
\[
0\to \mathcal A\to \Omega_X^1\to \omega_X\to \mathcal D\to 0.
\]
Since \(\bigwedge^2\mathbb{L}^{\bullet}_X\) has cohomology sheaves only in degrees \(-1\) and \(0\), its canonical truncation triangle is
\[
\mathcal A[1]\longrightarrow \bigwedge^2\mathbb{L}^{\bullet}_X\longrightarrow \mathcal D\longrightarrow \mathcal A[2].
\]
Because \(\mathcal A\) is supported on the finite set \(\mathrm{Sing}(X)\), one has \(H^i(X,\mathcal A)=0\) for all \(i>0\). Taking hypercohomology gives
\[
0\longrightarrow H^0\!\left(X,\bigwedge^2\mathbb{L}^{\bullet}_X\right)\longrightarrow H^0(X,\mathcal D)\longrightarrow 0,
\]
hence an isomorphism
\[
H^0\!\left(X,\bigwedge^2\mathbb{L}^{\bullet}_X\right)\cong H^0(X,\mathcal D).
\]
Under the same local description, \(\bigwedge^3\mathbb{L}^{\bullet}_A\) is represented by \(K[-1]\). Therefore
\[
\mathcal H^{-1}\!\left(\bigwedge^3\mathbb{L}^{\bullet}_X\right)_x
\cong
H^{-1}(K[-1])
=
H^0(K)
\cong
\mathcal D_x.
\]
By construction, for each \(x\in \mathrm{Sing}(X)\) the stalk map
\[
\alpha_x:\mathcal A_x\longrightarrow \mathcal D_x
\]
is the composite of \(\mathcal H^{-1}(d_{\mathrm{dR}})_x\) with the canonical identification
\[
\mathcal H^{-1}\!\left(\bigwedge^3\mathbb{L}^{\bullet}_X\right)_x
\cong
\mathcal D_x,
\]
and hence is the restriction of the local de Rham differential
\[
(\Omega_X^1)_x\longrightarrow
\mathcal H^0\!\left(\bigwedge^2\mathbb{L}^{\bullet}_X\right)_x=\mathcal D_x.
\]
Therefore, after taking global sections and using the identification above, the restriction of \(d_1^{1,0}\) along \(H^0(X,\mathcal A)\hookrightarrow H^0(X,\Omega_X^1)\) is identified with \(H^0(\alpha)\).
\end{proof}

We now specialize to the quasihomogeneous case. The remaining local input needed for the global degeneration criterion is the following surjectivity statement.

\begin{lemma}\label{lem:plane-local-surj}
Let \(x\in \mathrm{Sing}(X)\), and use the standing notation above.
Assume that \(x\) is quasihomogeneous. Then the stalk map of Lemma~\ref{lem:plane-d10-comparison}
\[
\alpha_x:\mathcal A_x\longrightarrow \mathcal D_x
\]
is surjective.
\end{lemma}

\begin{proof}
Since \(x\) is quasihomogeneous, we may choose a weighted homogeneous equation for \(x\), so that \eqref{eq:weighted-euler} holds.
By Saito's criterion~\eqref{eq:tau-mu-qh}, quasihomogeneity implies \(\tau_x=\mu_x\). Hence the natural quotient
\[
M_f\twoheadrightarrow T_f=M_f/fM_f
\]
is an isomorphism, since both sides have dimension \(\mu_x\).

Let \(\alpha_x:\mathcal A_x\to \mathcal D_x\) be the stalk at \(x\) of the morphism \(\alpha\) of Lemma~\ref{lem:plane-d10-comparison}. As in Lemmas~\ref{lem:plane-cotangent} and \ref{lem:plane-exterior-powers}, flat base change to the completion \(A=\widehat{\mathcal O}_{X,x}\) gives
\[
\mathcal A_x\otimes_{\mathcal O_{X,x}} A
\cong
H^{-1}\!\left(\bigwedge^2\mathbb{L}^{\bullet}_A\right),
\qquad
\mathcal D_x\otimes_{\mathcal O_{X,x}} A
\cong
H^{0}\!\left(\bigwedge^2\mathbb{L}^{\bullet}_A\right).
\]
Under the local model
\[
\bigwedge^2\mathbb{L}^{\bullet}_A\simeq K,
\]
the de Rham differential is represented by the chain map \(K\to K[-1]\) used in the proof of Lemma~\ref{lem:tail-map-euler}, hence \(\alpha_x\otimes_{\mathcal O_{X,x}} A\) identifies with the induced map
\[
H^{-1}(K)\longrightarrow H^0(K).
\]
By Lemma~\ref{lem:tail-map-euler}, this map is an isomorphism, hence surjective. Since \(A\) is faithfully flat over \(\mathcal O_{X,x}\), exactness of tensoring with \(A\) gives
\[
\operatorname{coker}(\alpha_x)\otimes_{\mathcal O_{X,x}} A
\cong
\operatorname{coker}(\alpha_x\otimes_{\mathcal O_{X,x}} A)=0.
\]
Therefore \(\operatorname{coker}(\alpha_x)=0\), and \(\alpha_x\) is surjective.
\end{proof}

Combining Lemma~\ref{lem:plane-d10-comparison}, Lemma~\ref{lem:plane-local-surj}, and the tail isomorphisms of Lemma~\ref{lem:tail-map-euler}, we obtain the following characterization of \(E_2\)-degeneration in the plane case.

\begin{proposition}\label{prop:plane-case}
Assume that \(X\) has only plane curve singularities. Then the Hodge-to-de Rham spectral sequence
\[
E_1^{p,q}(X)=H^q\!\left(X,\bigwedge^p\mathbb{L}^{\bullet}_X\right)
\Longrightarrow H^{p+q}_{\mathrm{sing}}(X^{\mathrm{an}},\mathbb C)
\]
degenerates at the \(E_2\)-page if and only if every singularity of $X$ is quasihomogeneous.
\end{proposition}

\begin{proof}
By Corollary~\ref{cor:plane-e1-shape}, the only possible nonzero differentials on the \(E_1\)-page are the maps \(d_1^{0,1}\), \(d_1^{1,0}\), and the negative-row tail maps
\[
d_1:H^{-p}\!\left(X,\bigwedge^{p+1}\mathbb{L}^{\bullet}_X\right)
\longrightarrow
H^{-p}\!\left(X,\bigwedge^{p+2}\mathbb{L}^{\bullet}_X\right),
\qquad p\ge 1.
\]
Let \(\mathcal A\), \(\mathcal D\), and \(\alpha:\mathcal A\to \mathcal D\) be as in Lemma~\ref{lem:plane-d10-comparison}. In particular, there is an exact sequence
\[
0\to \mathcal A\to \Omega_X^1\to \omega_X\to \mathcal D\to 0,
\]
and under the induced identification
\[
H^0\!\left(X,\bigwedge^2\mathbb{L}^{\bullet}_X\right)\cong H^0(X,\mathcal D),
\]
the restriction of \(d_1^{1,0}\) along \(H^0(X,\mathcal A)\hookrightarrow H^0(X,\Omega_X^1)\) is \(H^0(\alpha)\).
\begin{enumerate}
\item
Assume that every singularity is quasihomogeneous. At each singular point we may choose a weighted homogeneous local equation, so Lemma~\ref{lem:tail-map-euler} shows that every local tail map is an isomorphism. Since the negative-row groups are direct sums of these local contributions, all global tail maps are isomorphisms.
By Lemma~\ref{lem:plane-local-surj}, the map
\[
\alpha_x:\mathcal A_x\longrightarrow \mathcal D_x
\]
is surjective for every \(x\in \mathrm{Sing}(X)\). Since \(\mathcal A\) and \(\mathcal D\) are supported on the finite set \(\mathrm{Sing}(X)\), taking global sections gives a surjection
\[
H^0(X,\mathcal A)\twoheadrightarrow H^0(X,\mathcal D).
\]
By Lemma~\ref{lem:plane-d10-comparison}, this surjection is identified with the restriction of \(d_1^{1,0}\). Therefore \(d_1^{1,0}\) is surjective, so
\[
E_2^{2,0}(X)=0.
\]
Since the tail maps are isomorphisms, all negative-row terms also vanish on the \(E_2\)-page. Therefore the only possible nonzero \(E_2\)-terms are
\[
E_2^{0,0}(X),\qquad
E_2^{1,0}(X),\qquad
E_2^{0,1}(X),\qquad
E_2^{1,1}(X).
\]
Explicitly, the \(E_2\)-page is supported in the positions \((0,0)\), \((1,0)\), \((0,1)\), and \((1,1)\), and has the form
\[
\begin{tikzcd}[row sep=tiny, column sep=small]
\scriptstyle q=1 & E_2^{0,1}(X) & E_2^{1,1}(X) & 0
\\
\scriptstyle q=0 & E_2^{0,0}(X) & E_2^{1,0}(X) & 0
\\
\scriptstyle q=-1 & 0 & 0 & 0
\end{tikzcd}
\]
Any differential \(d_r\) with \(r\ge 2\) has bidegree \((r,1-r)\), so no such differential can start or end at one of these four positions. Hence \(E_2^{p,q}(X)=E_\infty^{p,q}(X)\) for all \((p,q)\).

\item
Conversely, assume that the spectral sequence degenerates at the \(E_2\)-page. Write
\[
\delta:=\sum_{i=1}^s\delta_{x_i},
\qquad
\tau:=\sum_{i=1}^s\tau_{x_i},
\qquad
R:=\sum_{i=1}^s(r_{x_i}-1).
\]
By Lemma~\ref{lem:tail-map-quotient}, for every \(p\ge 1\) the source and target of the global tail map have the same dimension \(\tau\). Moreover, in the local description of Lemma~\ref{lem:tail-map-quotient}, after forgetting the basis elements \(\gamma_i\) each complex \(K_p\) is the same Koszul complex \(K\), and the de Rham differential is induced by the same chain map \(K[-p+1]\to K[-p]\). Thus the local tail map is independent of \(p\), and hence so is the global tail map. If one of these tail maps were not an isomorphism, then none of them would be. Since the source and target have the same dimension, for every \(p\) this would leave a nonzero kernel at \((p+1,-p)\) and a nonzero cokernel at \((p+2,-p)\) on the \(E_2\)-page. Under the assumption \(E_2^{p,q}(X)=E_\infty^{p,q}(X)\) for all \((p,q)\), this would give infinitely many nonzero graded pieces of \(H^1_{\mathrm{sing}}(X^{\mathrm{an}},\mathbb C)\) and \(H^2_{\mathrm{sing}}(X^{\mathrm{an}},\mathbb C)\), which is a contradiction. Therefore all tail maps are isomorphisms, so every negative row vanishes on the \(E_2\)-page.

Next, we denote 
\[
u:=d_1^{1,0}:H^0(X,\Omega_X^1)\to H^0\!\left(X,\bigwedge^2\mathbb{L}^{\bullet}_X\right),
\]
\[
v:=d_1^{0,1}:H^1(X,\mathcal O_X)\to H^1(X,\Omega_X^1),
\]
and write
\[
\kappa:=\dim \ker(u),
\qquad
c:=\dim \operatorname{coker}(u).
\]
Let \(\nu:\widetilde X\to X\) be the normalization. Topologically, \(X^{\mathrm{an}}\) is obtained from \(\widetilde X^{\mathrm{an}}\) by identifying, for each \(i\), the \(r_{x_i}\) points of \(\nu^{-1}(x_i)\). Equivalently, \(X^{\mathrm{an}}\) is homotopy equivalent to \(\widetilde X^{\mathrm{an}}\) with \(R=\sum_i(r_{x_i}-1)\) additional circles attached. Since \(\widetilde X^{\mathrm{an}}\) is a compact Riemann surface of genus \(g\), it follows that
\[
b_1(X^{\mathrm{an}})=2g+R,
\qquad
b_2(X^{\mathrm{an}})=1.
\]
Indeed, choose pairwise disjoint contractible neighborhoods \(U_i\) of the singular points such that \(\nu^{-1}(U_i)\) is a disjoint union of \(r_{x_i}\) discs. Collapsing one disc in each fiber to the image point shows that, at \(x_i\), every identification after the first contributes one additional circle.
A Mayer--Vietoris computation then gives the displayed Betti numbers.
Therefore
\[
\dim H^1_{\mathrm{sing}}(X^{\mathrm{an}},\mathbb C)=2g+R,
\qquad
\dim H^2_{\mathrm{sing}}(X^{\mathrm{an}},\mathbb C)=1.
\]
Because the negative rows vanish on the \(E_2\)-page, the only terms contributing to total degree \(1\) are \(E_2^{0,1}(X)\) and \(E_2^{1,0}(X)\), while the only terms contributing to total degree \(2\) are \(E_2^{1,1}(X)\) and \(E_2^{2,0}(X)\). Hence
\[
\dim \ker(v)+\kappa=2g+R,
\qquad
\dim \operatorname{coker}(v)+c=1.
\]
Next we compute \(h^1(X,\Omega_X^1)\). From the exact sequence
\[
0\to \mathcal A\to \Omega_X^1\to \omega_X\to \mathcal D\to 0
\]
of Lemma~\ref{lem:plane-d10-comparison} and Lemma~\ref{lem:tail-map-quotient}, the first and last terms are supported on \(\mathrm{Sing}(X)\) and both have total length \(\tau\). Therefore
\[
\begin{aligned}
\chi(X,\Omega_X^1)
&=\chi(X,\omega_X)\\
&=h^0(X,\omega_X)-h^1(X,\omega_X)\\
&=h^1(X,\mathcal O_X)-h^0(X,\mathcal O_X)\\
&=p_a(X)-1\\
&=g+\delta-1,
\end{aligned}
\]
where the third equality is given by Serre duality.

Notice that
\[
H^0\!\left(X,\bigwedge^2\mathbb{L}^{\bullet}_X\right)=\bigoplus_i T_{x_i},
\]
by Lemma~\ref{lem:plane-d10-comparison} together with the local identification
\(\mathcal D_{x_i}\cong T_{x_i}\).
hence the target of \(u\) has dimension \(\tau\). Hence
\[
\dim \operatorname{im}(u)=\tau-c.
\]
By rank-nullity,
\[
h^0(X,\Omega_X^1)=\dim \ker(u)+\dim \operatorname{im}(u)=\kappa+\tau-c.
\]
Since
\[
\chi(X,\Omega_X^1)=h^0(X,\Omega_X^1)-h^1(X,\Omega_X^1),
\]
it follows that
\[
h^1(X,\Omega_X^1)=h^0(X,\Omega_X^1)-\chi(X,\Omega_X^1)
=\tau+1-g-\delta+\kappa-c.
\]

On the one hand,
\[
h^1(X,\mathcal O_X)=p_a(X)=g+\delta,
\]
so
\[
\operatorname{rank}(v)
=h^1(X,\mathcal O_X)-\dim \ker(v)
=(g+\delta)-(2g+R-\kappa)
=\delta-g-R+\kappa.
\]
On the other hand,
\[
\begin{aligned}
\operatorname{rank}(v)
&=h^1(X,\Omega_X^1)-\dim \operatorname{coker}(v)\\
&=(\tau+1-g-\delta+\kappa-c)-(1-c)\\
&=\tau-g-\delta+\kappa.
\end{aligned}
\]
Comparing these two expressions gives
\[
\tau=2\delta-R.
\]
Using Milnor's formula \eqref{eq:milnor-formula}, we conclude that
\[
\sum_{i=1}^s\tau_{x_i}
=
\tau
=
2\delta-R
=
\sum_{i=1}^s\mu_{x_i}.
\]
Since \(\tau_{x_i}\le \mu_{x_i}\) for every plane curve singularity, equality of the sums forces \(\tau_{x_i}=\mu_{x_i}\) for every \(i\). By Saito's criterion~\eqref{eq:tau-mu-qh}, every singularity of \(X\) is quasihomogeneous.
\end{enumerate}
\end{proof}

\subsection{Non-planar lci singularities}\label{subsec:non-planar-case}

In this section, we will show that a non-planar lci singularity forces a nonzero \(E_2\)-term in total degree \(>2\). We start by analyzing the local tail maps.

\begin{lemma}\label{lem:nonplanar-obstruction}
Let $A$ be a one-dimensional complete reduced local $\mathbb{C}$-algebra which is a local complete intersection, let \(\mathfrak{m}_A\) denote its unique maximal ideal. Assume that $A$ is not planar, and denote
\[
e:=\operatorname{embdim}(A)\ge 3.
\]
Then
\[
H^{-1}\!\left(\bigwedge^{e+1}\mathbb{L}^{\bullet}_A\right)\neq 0,
\]
and the de Rham differential induces a map on cohomology
\[
d_1:H^{-1}\!\left(\bigwedge^e\mathbb{L}^{\bullet}_A\right)\longrightarrow
H^{-1}\!\left(\bigwedge^{e+1}\mathbb{L}^{\bullet}_A\right),
\]
which is not surjective.
\end{lemma}

\begin{proof}
Since \(A\) is a one-dimensional lci singularity of embedding dimension \(e\), there exists a minimal Cohen presentation
\[
A\cong Q/(f_1,\dots,f_{e-1}),
\qquad
Q=\mathbb{C}[[x_1,\dots,x_e]],
\]
Its cotangent complex is represented by a minimal two-term complex
\[
\mathbb{L}^{\bullet}_A\simeq [F\xrightarrow{\partial} G],
\qquad
F=A^{e-1},\ G=A^e,
\]
where \(\partial\) is induced by the Jacobian matrix \(\left(\partial f_i/\partial x_j\right)\).
Since the Cohen presentation is minimal, the ideal \((f_1,\dots,f_{e-1})\) is contained in \(\mathfrak m_Q^2\). Hence each \(\partial f_i/\partial x_j\) lies in \(\mathfrak m_Q\), and its image in \(A\) lies in \(\mathfrak m_A\).

For \(p\ge 1\), write
\[
K_p:=\bigwedge^p\mathbb{L}^{\bullet}_A.
\]
Then \(K_p\) is represented by the standard complex for derived exterior powers of a two-term complex:
\[
\operatorname{Sym}^pF
\longrightarrow
\operatorname{Sym}^{p-1}F\otimes G
\longrightarrow
\operatorname{Sym}^{p-2}F\otimes \bigwedge^2 G
\longrightarrow \cdots \longrightarrow
\operatorname{Sym}^{p-e}F\otimes \bigwedge^e G,
\]
placed in degrees \([-p,-p+e]\), where terms with negative symmetric power are omitted. In particular, \(K_{e+1}\) ends with
\[
\operatorname{Sym}^2F\otimes \bigwedge^{e-1}G
\longrightarrow
F\otimes \bigwedge^e G
\longrightarrow 0,
\]
with \(F\otimes \bigwedge^e G\) in degree \(-1\). Since \(\bigwedge^{e+1}G=0\), there is no term in degree \(0\). Let
\[
\phi:\operatorname{Sym}^2F\otimes \bigwedge^{e-1}G\longrightarrow F\otimes \bigwedge^e G
\]
denote this last differential. Since this differential is functorially induced by \(\partial\), all entries of \(\phi\) lie in \(\mathfrak m_A\). Hence
\[
\operatorname{im}(\phi)\subseteq \mathfrak m_A\!\left(F\otimes \bigwedge^e G\right).
\]
Because \(K_{e+1}\) has no term in degree \(0\), one has
\[
H^{-1}(K_{e+1})=\operatorname{coker}(\phi).
\]
Because \(F\otimes \bigwedge^e G\neq 0\), Nakayama's lemma gives
\[
H^{-1}(K_{e+1})
=
\operatorname{coker}(\phi)\neq 0.
\]

Now let \(k:=A/\mathfrak m_A\), and consider the de Rham chain map
\[
\delta_e:K_e\longrightarrow K_{e+1}.
\]
Using the standard Tate dg-resolution
\[
S=Q[\epsilon_1,\dots,\epsilon_{e-1}],
\qquad
\partial(\epsilon_i)=f_i,
\]
one has \(\mathbb{L}^{\bullet}_A\simeq \Omega^1_{S}\otimes_S A\), where \(\Omega^1_{S}\) is free on the generators \(d\epsilon_i\) and \(dx_j\). The complexes representing \(K_e\) and \(K_{e+1}\) are functorially induced from this two-term dg-model, so each term is free with a basis consisting of symmetric monomials in the \(d\epsilon_i\) and wedge monomials in the \(dx_j\). Since
\[
d_{\mathrm{dR}}(d\epsilon_i)=d_{\mathrm{dR}}(dx_j)=0,
\]
every such basis monomial is \(d_{\mathrm{dR}}\)-closed. By functoriality, the induced chain map \(\delta_e\) is therefore obtained by applying the ordinary de Rham differential only to the coefficients in \(A\). After tensoring with \(k\), all coefficients become constants, so every coefficient differential vanishes. Hence
\[
\delta_e\otimes_A \mathrm{id}_k:
K_e\otimes_A k\longrightarrow K_{e+1}\otimes_A k
\]
is the zero map.

Since \(K_{e+1}\) is a bounded complex of free \(A\)-modules, the universal coefficient short exact sequence in degree \(-1\) gives
\[
0\longrightarrow H^{-1}(K_{e+1})\otimes_A k
\xrightarrow{\iota}
H^{-1}(K_{e+1}\otimes_A k)
\longrightarrow
\operatorname{Tor}_1^A(H^0(K_{e+1}),k)
\longrightarrow 0.
\]
Since \(K_{e+1}\) has no term in degree \(0\), one has \(H^0(K_{e+1})=0\). Hence the \(\operatorname{Tor}_1\)-term vanishes, and \(\iota\) is injective.

Next, we consider the following commutative diagram:
\[
\begin{tikzcd}
H^{-1}(K_e) \ar[r, "d_1"] \ar[d, "\rho_e"']
& H^{-1}(K_{e+1}) \ar[d, "\rho_{e+1}"] \\
H^{-1}(K_e\otimes_A k) \ar[r, "0"']
& H^{-1}(K_{e+1}\otimes_A k)
\end{tikzcd}
\]
where the bottom map is induced from the map \(\delta_e\otimes_A k\), hence is the zero map. Let
\[
q:H^{-1}(K_{e+1})\twoheadrightarrow H^{-1}(K_{e+1})\otimes_A k,
\]
be the quotient map. Then \(\rho_{e+1}=\iota\circ q\). From the commutative square, one has \(\rho_{e+1}\circ d_1=0\). Hence
\[
0=\rho_{e+1}\circ d_1=\iota\circ q\circ d_1.
\]
Because \(\iota\) is injective, it follows that \(q\circ d_1=0\), i.e.
\[
\operatorname{im}(d_1)\subseteq \mathfrak m_A H^{-1}(K_{e+1}).
\]
Since \(H^{-1}(K_{e+1})\neq 0\), Nakayama's lemma implies that \(d_1\) is not surjective.
\end{proof}

\begin{corollary}\label{cor:nonplanar-global-obstruction}
Let $X$ be an integral projective curve over $\mathbb{C}$ with a non-planar lci singularity. Then the Hodge-to-de Rham spectral sequence of $X$ does not degenerate at the \(E_2\)-page.
\end{corollary}

\begin{proof}
Set
\[
\displaystyle
e:=\max_{x\in \mathrm{Sing}(X)} \operatorname{embdim}\!\left(\widehat{\mathcal O}_{X,x}\right).
\]
Since \(X\) has a non-planar singularity, one has \(e\ge 3\).

For \(p\ge 2\), the same support-on-the-singular-locus argument as in Lemma~\ref{lem:plane-exterior-powers} shows that \(\bigwedge^p\mathbb{L}^{\bullet}_X\) is supported on the finite singular set. Hence
\[
E_1^{p,-1}(X)
=
H^{-1}\!\left(X,\bigwedge^p\mathbb{L}^{\bullet}_X\right)
\cong
H^0\!\left(X,\mathcal H^{-1}\!\left(\bigwedge^p\mathbb{L}^{\bullet}_X\right)\right)
\]
and the global \(d_1\)-map on this row is induced by the corresponding morphism of cohomology sheaves.

Fix a singular point \(x\in \mathrm{Sing}(X)\), and write
\[
\widehat{\mathcal O}_{X,x}\cong Q_x/(f_1,\dots,f_{e_x-1}),
\qquad
e_x:=\operatorname{embdim}\!\left(\widehat{\mathcal O}_{X,x}\right)\le e.
\]
Then the completed stalk of \(\bigwedge^p\mathbb{L}^{\bullet}_X\) at \(x\) is represented by the corresponding complex \(K_p(x)\), which is concentrated in degrees \([-p,-p+e_x]\). Therefore
\[
H^{-1}\!\left(K_p(x)\right)=0
\qquad\text{whenever}\qquad
p\ge e_x+2.
\]
Applying this with \(p=e+2\), we obtain
\[
\mathcal H^{-1}\!\left(\bigwedge^{e+2}\mathbb{L}^{\bullet}_X\right)_x=0
\qquad\text{for every }x\in \mathrm{Sing}(X),
\]
and hence
\[
E_1^{e+2,-1}(X)=0.
\]

Similarly, if \(e_x<e\), then \(H^{-1}(K_{e+1}(x))=0\). Choose a singular point \(x_0\) with \(e_{x_0}=e\). By Lemma~\ref{lem:nonplanar-obstruction},
\[
H^{-1}\!\left(K_{e+1}(x_0)\right)\neq 0,
\]
so faithful flatness of completion gives
\[
\mathcal H^{-1}\!\left(\bigwedge^{e+1}\mathbb{L}^{\bullet}_X\right)_{x_0}\neq 0.
\]
Therefore
\[
E_1^{e+1,-1}(X)\neq 0.
\]

Thus the relevant part of the \(q=-1\) row of the \(E_1\)-page has the form
\[
\begin{tikzcd}[row sep=tiny, column sep=small]
& \scriptstyle p=e & \scriptstyle p=e+1 & \scriptstyle p=e+2 \\
\scriptstyle q=0
& E_1^{e,0}(X) \ar[r]
& E_1^{e+1,0}(X) \ar[r]
& E_1^{e+2,0}(X)
\\
\scriptstyle q=-1
& E_1^{e,-1}(X) \ar[r,"u"]
& E_1^{e+1,-1}(X) \ar[r]
& 0
\end{tikzcd}
\]
with \(E_1^{e+1,-1}(X)\neq 0\). It remains only to show that
\(u\) is not surjective.
Because \(\mathcal H^{-1}(\bigwedge^e\mathbb{L}^{\bullet}_X)\) and \(\mathcal H^{-1}(\bigwedge^{e+1}\mathbb{L}^{\bullet}_X)\) are supported on the finite set \(\mathrm{Sing}(X)\), taking global sections identifies \(u\) with the direct sum of the induced stalk maps over the singular points. Projecting the target onto the \(x_0\)-summand, we obtain a component map which becomes, after completion,
\[
d_1:H^{-1}\!\left(K_e(x_0)\right)\longrightarrow H^{-1}\!\left(K_{e+1}(x_0)\right),
\]
which is not surjective by Lemma~\ref{lem:nonplanar-obstruction}. Therefore \(u\) itself cannot be surjective. So
\[
E_2^{e+1,-1}(X)
=
\operatorname{coker}(u)\neq 0.
\]

This term has total degree \(e\). Since \(e\ge 3\) and \(X\) is a projective curve, one has
\[
H^e_{\mathrm{sing}}(X^{\mathrm{an}},\mathbb C)=0.
\]
Therefore \(E_\infty^{e+1,-1}(X)=0\), so the spectral sequence cannot satisfy \(E_2^{p,q}(X)=E_\infty^{p,q}(X)\) for all \((p,q)\). Hence it does not degenerate at the \(E_2\)-page.
\end{proof}

\subsection{Main theorem}\label{subsec:main-theorem}

We can now combine the plane case with the non-planar obstruction.

\begin{theorem}\label{thm:main}
For an integral projective curve \(X\) with local complete intersection singularities, its Hodge-to-de Rham spectral sequence
\[
E_1^{p,q}(X)=H^q\!\left(X,\bigwedge^p\mathbb{L}^{\bullet}_X\right)
\Longrightarrow H^{p+q}_{\mathrm{sing}}(X^{\mathrm{an}},\mathbb C)
\]
degenerates at the \(E_2\)-page if and only if every singularity of \(X\) is a quasihomogeneous plane curve singularity.
\end{theorem}

\begin{proof}
If every singularity of \(X\) is a quasihomogeneous plane curve singularity, then the claim follows from Proposition~\ref{prop:plane-case}.

Conversely, assume that the Hodge-to-de Rham spectral sequence of \(X\) degenerates at the \(E_2\)-page. Since all singularities of \(X\) are local complete intersections by hypothesis, Corollary~\ref{cor:nonplanar-global-obstruction} excludes any non-planar singularity. Therefore all singularities of \(X\) are planar, and Proposition~\ref{prop:plane-case} implies that they are quasihomogeneous.
\end{proof}


\subsection{The Hochschild-to-cyclic spectral sequence}\label{subsec:HC}

In this section, we study the degeneration of the Hochschild-to-cyclic spectral sequence and prove an analogous criterion for degeneration at the \(E_2\)-page.

We start by briefly recalling the \emph{Hochschild-to-cyclic spectral sequence} in the form needed below.

A mixed complex in an abelian category \(\cE\) is a complex \((V_{\bullet},b)\) together with a
morphism of complexes \(B:V_{\bullet}\to V_{\bullet}[-1]\) satisfying \(B^2=0\); see
\cite[\S~1]{Kassel_87} and \cite[\S~1.1]{Kaledin_17}. A basic example is the Hochschild chain
complex \(\mathrm{CC}_{\bullet}(A)=A^{\otimes(\bullet+1)}\) of an algebra \(A\), equipped with the
Hochschild differential \(b\) and Connes' operator \(B\); see \cite{Connes_86} and
\cite[\S~2]{Kassel_87}.\footnote{It is denoted as the \((b,B)\)-bicomplex \(\mathcal{B}(A)\) in \cite[\S~2]{Loday_92}.} 
Keller globalized this construction to a mixed complex \((\mathrm{CC}_{\bullet}(W),b,B)\) for any variety \(W\) \cite[\S~2.1, \S~5.2]{Keller_98}.\footnote{In modern language, the sheaf Hochschild complex on \(W\) is
\[
\underline{\mathsf{HH}}_W:=\mathcal{O}_W\otimes^{\mathbf L}_{\mathcal{O}_{W\times W}}\mathcal{O}_W,
\]
and Keller's mixed complex \((\mathrm{CC}_{\bullet}(W),b,B)\) is a global model computing
\(R\Gamma(W,\underline{\mathsf{HH}}_W)\).}

Let \(u\) be a formal variable of homological degree \(-2\). For any mixed complex
\((V_{\bullet},b,B)\), the \emph{negative cyclic complex} is the total complex
\[
\bigl(V_{\bullet}[[u]],\,b+uB\bigr),
\]
and its homology is called the \emph{negative cyclic homology} \(\HC^-_*(V_{\bullet})\). The \(u\)-adic
filtration induces a convergent spectral sequence \cite[\S~1.2]{Kaledin_17}
\[
E_{1,a,b}^{\mathrm{HC}}(V_\bullet)=u^aH_{a+b}(V_{\bullet})\Longrightarrow \HC^-_{a+b}(V_{\bullet}).\footnote{Equivalently, one may first take homology with respect to the vertical differential \(b\)
for the double complex \(V_{\bullet}[[u]]\); this gives the \(E_1\)-page written below and converges to \(\HC^-_*(V_{\bullet})\).}
\]

For a variety \(W\), applying this construction to \(\mathrm{CC}_{\bullet}(W)\) gives the
\emph{Hochschild-to-cyclic spectral sequence}
\[
E_{1,a,b}^{\mathrm{HC}}(W)=u^a\HH_{a+b}(W)\Longrightarrow \HC^-_{a+b}(W).
\]
We write
\[
E_{r,a,b}^{\mathrm{HC}}(X)
\]
for the \(r\)-th page of the Hochschild-to-cyclic spectral sequence of \(X\). In particular, its \(E_1\)-page has the following form:
\[
\begin{tikzcd}[row sep=tiny]
& \vdots & \vdots & \vdots \\
\cdots & 0 \ar[r] & \HH_{-1} \ar[r,"uB"] & u\,\HH_{0} \ar[r,"uB"] & \cdots \\
\cdots & 0 \ar[r] & \HH_{0} \ar[r,"uB"] & u\,\HH_{1} \ar[r,"uB"] & \cdots \\
\cdots & 0 \ar[r] & \HH_{1} \ar[r,"uB"] & u\,\HH_{2} \ar[r,"uB"] & \cdots \\
& \vdots & \vdots & \vdots
\end{tikzcd}
\]

\begin{proposition}\label{prop:HC-to-HdR}
Let \(X\) be a separated scheme of finite type over \(\mathbb C\). If its Hochschild-to-cyclic spectral sequence
degenerates at the
\(E_k\)-page, then the Hodge-to-de Rham spectral sequence of \(X\) also degenerates
at the \(E_k\)-page.
\end{proposition}

\begin{proof}
The point is to compare the two spectral sequences at the level of filtered complexes, not only at
the level of their \(E_1\)-pages.

Let
\[
\mathsf{C}_{\mathrm{Hoch}}(X):=\bigl(\mathrm{CC}_{\bullet}(X),b,B\bigr)
\]
be Keller's Hochschild mixed complex of \(X\), and let
\[
\underline{\mathsf{HH}}_X:=\mathcal{O}_X\otimes^L_{\mathcal{O}_{X\times X}}\mathcal{O}_X
\]
be the sheaf Hochschild complex on \(X\). By the discussion above,
\(\mathsf{C}_{\mathrm{Hoch}}(X)\) computes \(R\Gamma(X,\underline{\mathsf{HH}}_X)\).
Since \(X\) is separated, it is in particular semi-separated. Therefore
To\"en--Vezzosi \cite[Corollary~1.2 and the discussion following it]{Toen-Vezzosi_11}
give a natural multiplicative HKR-type comparison at the level of sheaves:
there is an isomorphism in the homotopy category of sheaves of commutative
\(\mathcal{O}_X\)-dg-algebras
\[
\operatorname{Sym}^{\bullet}_{\mathcal{O}_X}\!\left(\mathbb{L}^{\bullet}_X[1]\right)
\simeq
\underline{\mathsf{HH}}_X,
\]
and the discussion following Corollary~1.2 shows that this comparison is compatible with the
\(S^1\)-action on the Hochschild side and with the de Rham differential. Let
\[
\mathcal{E}_X:=
\left(
\bigoplus_{p\ge 0}\bigwedge^p\mathbb{L}^{\bullet}_X[-p],
d_{\mathrm{int}},
d_{\mathrm{dR}}
\right)
\]
denote the derived de Rham mixed complex on \(X\), equivalently
\[
\mathcal{E}_X\cong \left(\operatorname{Sym}^{\bullet}_{\mathcal{O}_X}\!\left(\mathbb{L}^{\bullet}_X[1]\right),d_{\mathrm{int}},d_{\mathrm{dR}}\right).
\]
Forgetting the dg-algebra structures, this gives an isomorphism in the homotopy category of sheaves of mixed complexes
\[
\mathcal{E}_X\simeq \underline{\mathsf{HH}}_X.
\]
Applying \(R\Gamma(X,-)\), we obtain an isomorphism in the derived category of mixed complexes
\[
R\Gamma(X,\mathcal{E}_X)\simeq R\Gamma(X,\underline{\mathsf{HH}}_X).
\]
Since Keller's mixed complex \(\mathsf{C}_{\mathrm{Hoch}}(X)\) is a model for the right-hand side,
we may replace it by the quasi-isomorphic mixed complex
\[
\mathsf{E}(X):=R\Gamma(X,\mathcal{E}_X)=
\left(
\bigoplus_{p\ge 0}R\Gamma\!\left(X,\bigwedge^p\mathbb{L}^{\bullet}_X\right)[-p],
d_{\mathrm{int}},
d_{\mathrm{dR}}
\right),
\]
where \(d_{\mathrm{int}}\) is induced by the internal differential on
\(\bigwedge^p\mathbb{L}^{\bullet}_X[-p]\), and \(d_{\mathrm{dR}}\) is induced by the de Rham differential.
Choose a finite affine open cover with affine intersections. Since derived global sections of quasi-coherent complexes are computed by the corresponding finite \v{C}ech complex, and this finite \v{C}ech complex commutes with arbitrary direct sums, the displayed decomposition follows.

Applying the negative cyclic construction recalled above to \(\mathsf{E}(X)\) yields the filtered complex
\[
\left(
\prod_{a\ge 0}u^a\bigoplus_{p\ge 0}R\Gamma\!\left(X,\bigwedge^p\mathbb{L}^{\bullet}_X\right)[-p],
d_{\mathrm{int}}+u\,d_{\mathrm{dR}}
\right),
\]
with filtration given by powers of \(u\). This construction is functorial in mixed complexes, and a quasi-isomorphism of mixed complexes induces a filtered quasi-isomorphism on the resulting negative cyclic complexes: indeed, the associated graded for the \(u\)-adic filtration is \(\bigoplus_{a\ge 0}u^a(V_{\bullet},b)\). Therefore the spectral sequence associated with this filtered complex is the Hochschild-to-cyclic spectral sequence of \(X\).


Now set
\[
m:=p-a.
\]
The differential \(d_{\mathrm{int}}\) preserves both \(a\) and \(p\), while \(u\,d_{\mathrm{dR}}\) sends the summand with indices \((a,p)\) to the summand with indices \((a+1,p+1)\). Hence both preserve \(m\), so this filtered complex splits as a product of filtered subcomplexes indexed by \(m\). This decomposition is compatible with the \(u\)-adic filtration, since each filtration piece is the corresponding product of the filtration pieces on the \(m\)-summands. In particular, projection to the summands with \(p=a\) defines a filtered direct summand
\[
\left(
\prod_{i\ge 0}u^iR\Gamma\!\left(X,\bigwedge^i\mathbb{L}^{\bullet}_X\right)[-i],
d_{\mathrm{int}}+u\,d_{\mathrm{dR}}
\right).
\]
After forgetting the formal symbols \(u^i\), this is exactly
\[
R\Gamma\!\left(X,\widehat{\mathrm{dR}}_X^{\bullet}\right)
=
\prod_{i\ge 0}R\Gamma\!\left(X,\bigwedge^i\mathbb{L}^{\bullet}_X\right)[-i]
\]
with its Hodge filtration. Therefore the spectral sequence of the \(m=0\) summand is precisely the Hodge-to-de Rham spectral sequence of \(X\).

Since this is a filtered direct summand of the filtered complex computing the
Hochschild-to-cyclic spectral sequence, every page of its spectral sequence is a direct summand of the corresponding page of the Hochschild-to-cyclic spectral sequence. Hence if the latter degenerates at the \(E_k\)-page, then so does the former. Therefore the Hodge-to-de Rham spectral sequence of \(X\) degenerates at the \(E_k\)-page.
\end{proof}

\begin{remark}\label{rem:HC-filtered-complex}
The use of the mixed complex \(\mathsf{C}_{\mathrm{Hoch}}(X)\) is essential here. Knowing only the \(E_1\)-page and its differential determines the \(E_2\)-page,
but does not in general control the higher differentials.
\end{remark}

\begin{proposition}\label{prop:qhom-implies-HC}
Let \(X\) be an integral projective curve over \(\mathbb C\) with local complete intersection singularities. Assume that every singularity of \(X\) is a quasihomogeneous plane curve singularity. Then the Hochschild-to-cyclic spectral sequence of \(X\) degenerates at the \(E_2\)-page.
\end{proposition}

\begin{proof}
By the proof of Proposition~\ref{prop:HC-to-HdR}, the Hochschild-to-cyclic spectral sequence of \(X\) is computed by the filtered complex
\[
\left(
\prod_{a\ge 0}u^a\bigoplus_{p\ge 0}R\Gamma\!\left(X,\bigwedge^p\mathbb{L}^{\bullet}_X\right)[-p],
d_{\mathrm{int}}+u\,d_{\mathrm{dR}}
\right),
\]
and this filtered complex splits according to the integer
\[
m:=p-a.
\]
For each \(m\in \mathbb Z\), let \(\mathsf F_m(X)\) denote the filtered subcomplex formed by the summands with \(p-a=m\):
\[
\mathsf F_m(X):=
\left(
\prod_{a\ge \max(0,-m)}u^aR\Gamma\!\left(X,\bigwedge^{a+m}\mathbb{L}^{\bullet}_X\right)[-a-m],
d_{\mathrm{int}}+u\,d_{\mathrm{dR}}
\right).
\]
Since the differential preserves \(m\), it is enough to prove that the spectral sequence of each \(\mathsf F_m(X)\) degenerates at the \(E_2\)-page.

Taking homology with respect to \(d_{\mathrm{int}}\), the \(E_1\)-page of
\(\mathsf F_m(X)\) is obtained from the Hodge-to-de Rham \(E_1\)-page by keeping
only the columns with \(p\ge \max(0,m)\) and placing the \(p\)-th column in filtration degree \(a=p-m\). The \(d_1\)-differential is induced by \(u\,d_{\mathrm{dR}}\), hence by the ordinary de Rham differential on these groups.
Equivalently, for fixed \(m\), the filtration degree \(a\) of
\(E_1\bigl(\mathsf F_m(X)\bigr)\) is the Hodge column \(p=a+m\), and the only
nonzero \(E_1\)-differentials are the shifted copies
\[
H^q\!\left(X,\bigwedge^{a+m}\mathbb{L}^{\bullet}_X\right)\longrightarrow
H^q\!\left(X,\bigwedge^{a+m+1}\mathbb{L}^{\bullet}_X\right)
\]
of the \(d_1\)-maps in the Hodge-to-de Rham spectral sequence.

Since every singularity of \(X\) is quasihomogeneous and planar, Lemma~\ref{lem:tail-map-euler} shows that every local negative-row tail map is an isomorphism, hence all global tail maps are isomorphisms. Moreover, by Lemma~\ref{lem:plane-local-surj} and the global argument in the proof of Proposition~\ref{prop:plane-case}, the map
\[
H^0(X,\Omega_X^1)\longrightarrow H^0\!\left(X,\bigwedge^2\mathbb{L}^{\bullet}_X\right)
\]
is surjective.

\begin{itemize}
  \item If \(m\ge 3\), then every surviving column lies in the tail region \(p\ge 3\). Hence every \(E_1\)-differential on the spectral sequence of \(\mathsf F_m(X)\) is an isomorphism, so
  \[
    E_2\bigl(\mathsf F_m(X)\bigr)=0.
  \]
  \item If \(m=2\), then the only non-tail term is the isolated group
  \[
    H^0\!\left(X,\bigwedge^2\mathbb{L}^{\bullet}_X\right).
  \]
  All remaining columns lie in the tail region \(p\ge 3\), where the \(d_1\)-maps are isomorphisms. Therefore the \(E_2\)-page of \(\mathsf F_2(X)\) is supported in a single position, so again no higher differential can occur.
  \item If \(m=1\), then the \(p=0\) column is absent. The tail maps are still isomorphisms, and the only remaining \(d_1\)-map in nonnegative rows is the surjection
  \[
    H^0(X,\Omega_X^1)\longrightarrow H^0\!\left(X,\bigwedge^2\mathbb{L}^{\bullet}_X\right).
  \]
  Therefore the \(E_2\)-page of \(\mathsf F_1(X)\) is supported in only two positions, namely those coming from \(H^0(X,\Omega_X^1)\) and \(H^1(X,\Omega_X^1)\). No higher differential can start or end at either position.
  \item If \(m\le 0\), then all columns \(p\ge 0\) occur. Thus the \(E_1\)-page of
  \(\mathsf F_m(X)\) has exactly the same shape as in the first part of
  Proposition~\ref{prop:plane-case}, up to a horizontal shift by \(-m\) in the filtration direction. The same argument therefore shows that its \(E_2\)-page is supported in only four positions, and no higher differential can have source or target among them.
\end{itemize}
Hence every \(\mathsf F_m(X)\) degenerates at the \(E_2\)-page. Since the full Hochschild-to-cyclic filtered complex is the product of these \(\mathsf F_m(X)\), its spectral sequence also degenerates at the \(E_2\)-page: cycles and boundaries are computed componentwise in this product decomposition, so each page is the product of the corresponding pages of the spectral sequences of the \(\mathsf F_m(X)\).
\end{proof}

\begin{corollary}\label{cor:equiv-HdR-HC}
Let \(X\) be an integral projective curve over \(\mathbb C\) with local complete intersection
singularities. Then the following are equivalent:
\begin{enumerate}
\item the Hodge-to-de Rham spectral sequence of \(X\) degenerates at the \(E_2\)-page;
\item the Hochschild-to-cyclic spectral sequence of \(X\) degenerates at the \(E_2\)-page;
\item every singularity of \(X\) is a quasihomogeneous plane curve singularity.
\end{enumerate}
\end{corollary}

\begin{proof}
The equivalence of (1) and (3) is Theorem~\ref{thm:main}. 
The implication
\((2)\Longrightarrow (1)\)
is Proposition~\ref{prop:HC-to-HdR}, and the implication
\((3)\Longrightarrow (2)\)
is Proposition~\ref{prop:qhom-implies-HC}.
\end{proof}

\section{Further directions}\label{sec:further-directions}

The results of this paper suggest several natural questions beyond the case of integral projective lci curves. We record them here as possible further directions rather than precise conjectures.

\begin{enumerate}
\item \textbf{Nonreduced curves.}
Let $X$ be a projective lci curve over \(\mathbb C\) which is not necessarily reduced. Is there a singularity-theoretic criterion for the \(E_2\)-degeneration of the derived Hodge-to-de Rham spectral sequence of $X$? Even for very simple nonreduced curves, such as square-zero thickenings of reduced curves, the derived exterior powers of the cotangent complex can behave quite differently from the reduced case. It would therefore be interesting to understand whether the presence of nilpotents forces new differentials, and whether any analogue of Theorem~\ref{thm:main} survives in this setting.

\item \textbf{Higher-dimensional lci varieties.}
Let $X$ be a projective lci variety over \(\mathbb C\) of dimension at least $2$. Can one characterize the \(E_2\)-degeneration of the derived Hodge-to-de Rham spectral sequence of $X$ in terms of the singularities of $X$? At present, even the first nontrivial cases appear widely open. A reasonable place to start would be projective surfaces with isolated lci singularities, or more specifically isolated hypersurface singularities. One may also ask whether quasihomogeneity continues to play a distinguished role in higher dimensions.

\item \textbf{Degeneration at later pages.}
For singular projective varieties, can one determine the first page at which the derived Hodge-to-de Rham spectral sequence degenerates in terms of the local singularity types? Theorem~\ref{thm:main} gives a complete criterion for \(E_2\)-degeneration in the lci curve case. A natural next problem is to ask whether there exists a local invariant \(\Delta\) of a singularity that controls the first possible nonzero differential, or more generally the minimal page of degeneration. One may hope for a function \(r=r(n,\Delta)\), depending only on the dimension \(n\) and a suitable local invariant \(\Delta\), such that degeneration always occurs by the \(E_r\)-page.
\end{enumerate}

We hope that the techniques developed here for lci curves, together with the relation to the Hochschild-to-cyclic spectral sequence established in Subsection~\ref{subsec:HC}, may provide a starting point for some of these questions. 

\printbibliography

\end{document}
